\documentclass[10pt,a4paper]{amsart}
\usepackage{amsmath,amssymb,amsthm,graphicx}
\usepackage{verbatim}
\newenvironment{reference}{\comment}{\endcomment}
%\newenvironment{reference}{}{}
\newenvironment{slogan}{\comment}{\endcomment}
\newenvironment{history}{\comment}{\endcomment}
% For commutative diagrams we use Xy-pic
\usepackage[all]{xy}

% We use 2cell for 2-commutative diagrams.
\xyoption{2cell}
% Theorem environments.
%
\theoremstyle{plain}
\newtheorem{theorem}[subsection]{Theorem}
\newtheorem{proposition}[subsection]{Proposition}
\newtheorem{lemma}[subsection]{Lemma}

\theoremstyle{definition}
\newtheorem{definition}[subsection]{Definition}
\newtheorem{example}[subsection]{Example}
\newtheorem{exercise}[subsection]{Exercise}
\newtheorem{situation}[subsection]{Situation}

\theoremstyle{remark}
\newtheorem{remark}[subsection]{Remark}
\newtheorem{remarks}[subsection]{Remarks}

\numberwithin{equation}{subsection}

% Macros
%
\def\lim{\mathop{\mathrm{lim}}\nolimits}
\def\colim{\mathop{\mathrm{colim}}\nolimits}
\def\Spec{\mathop{\mathrm{Spec}}}
\def\Hom{\mathop{\mathrm{Hom}}\nolimits}
\def\Ext{\mathop{\mathrm{Ext}}\nolimits}
\def\SheafHom{\mathop{\mathcal{H}\!\mathit{om}}\nolimits}
\def\SheafExt{\mathop{\mathcal{E}\!\mathit{xt}}\nolimits}
\def\Sch{\mathit{Sch}}
\def\Mor{\mathop{\mathrm{Mor}}\nolimits}
\def\Ob{\mathop{\mathrm{Ob}}\nolimits}
\def\Sh{\mathop{\mathit{Sh}}\nolimits}
\def\NL{\mathop{N\!L}\nolimits}
\def\CH{\mathop{\mathrm{CH}}\nolimits}
\def\proetale{{pro\text{-}\acute{e}tale}}
\def\etale{{\acute{e}tale}}
\def\QCoh{\mathit{QCoh}}
\def\Ker{\mathop{\mathrm{Ker}}}
\def\Im{\mathop{\mathrm{Im}}}
\def\Coker{\mathop{\mathrm{Coker}}}
\def\Coim{\mathop{\mathrm{Coim}}}

% Boxtimes
%
\DeclareMathSymbol{\boxtimes}{\mathbin}{AMSa}{"02}

%
% Macros for moduli stacks/spaces
%
\def\QCohstack{\mathcal{QC}\!\mathit{oh}}
\def\Cohstack{\mathcal{C}\!\mathit{oh}}
\def\Spacesstack{\mathcal{S}\!\mathit{paces}}
\def\Quotfunctor{\mathrm{Quot}}
\def\Hilbfunctor{\mathrm{Hilb}}
\def\Curvesstack{\mathcal{C}\!\mathit{urves}}
\def\Polarizedstack{\mathcal{P}\!\mathit{olarized}}
\def\Complexesstack{\mathcal{C}\!\mathit{omplexes}}
% \Pic is the operator that assigns to X its picard group, usage \Pic(X)
% \Picardstack_{X/B} denotes the Picard stack of X over B
% \Picardfunctor_{X/B} denotes the Picard functor of X over B
\def\Pic{\mathop{\mathrm{Pic}}\nolimits}
\def\Picardstack{\mathcal{P}\!\mathit{ic}}
\def\Picardfunctor{\mathrm{Pic}}
\def\Deformationcategory{\mathcal{D}\!\mathit{ef}}
\usepackage{cleveref}
\crefname{originallemma}{original lemma}{original lemmas}
\Crefname{originallemma}{Original lemma}{Original lemmas}
\makeatletter
\newlabel{algebra-lemma-CM-over-regular-flat}{{00R4}{source}{}{sourcepointer-1}{}}
\newlabel{algebra-lemma-CM-over-regular-flat@cref}{{[originallemma][0][]00R4}{[1][0][]source}{}{}{}}
\makeatother
\crefname{originallemma}{original lemma}{original lemmas}
\Crefname{originallemma}{Original lemma}{Original lemmas}
\makeatletter
\newlabel{etale-lemma-relative-frobenius-etale}{{0EBS}{source}{}{sourcepointer-2}{}}
\newlabel{etale-lemma-relative-frobenius-etale@cref}{{[originallemma][0][]0EBS}{[1][0][]source}{}{}{}}
\makeatother
\crefname{originalreference}{original reference}{original references}
\Crefname{originalreference}{Original reference}{Original references}
\makeatletter
\newlabel{example-Garel}{{0FLD}{source}{}{sourcepointer-3}{}}
\newlabel{example-Garel@cref}{{[originalreference][0][]0FLD}{[1][0][]source}{}{}{}}
\makeatother
\crefname{originallemma}{original lemma}{original lemmas}
\Crefname{originallemma}{Original lemma}{Original lemmas}
\makeatletter
\newlabel{lemma-etale}{{0FLV}{source}{}{sourcepointer-4}{}}
\newlabel{lemma-etale@cref}{{[originallemma][0][]0FLV}{[1][0][]source}{}{}{}}
\makeatother
\crefname{originallemma}{original lemma}{original lemmas}
\Crefname{originallemma}{Original lemma}{Original lemmas}
\makeatletter
\newlabel{more-morphisms-lemma-flat-lci}{{069K}{source}{}{sourcepointer-5}{}}
\newlabel{more-morphisms-lemma-flat-lci@cref}{{[originallemma][0][]069K}{[1][0][]source}{}{}{}}
\makeatother
\crefname{originallemma}{original lemma}{original lemmas}
\Crefname{originallemma}{Original lemma}{Original lemmas}
\makeatletter
\newlabel{more-morphisms-lemma-lci-permanence}{{069M}{source}{}{sourcepointer-6}{}}
\newlabel{more-morphisms-lemma-lci-permanence@cref}{{[originallemma][0][]069M}{[1][0][]source}{}{}{}}
\makeatother
\crefname{originaltheorem}{original theorem}{original theorems}
\Crefname{originaltheorem}{Original theorem}{Original theorems}
\makeatletter
\newlabel{more-morphisms-theorem-criterion-flatness-fibre}{{039C}{source}{}{sourcepointer-7}{}}
\newlabel{more-morphisms-theorem-criterion-flatness-fibre@cref}{{[originaltheorem][0][]039C}{[1][0][]source}{}{}{}}
\makeatother
\crefname{originallemma}{original lemma}{original lemmas}
\Crefname{originallemma}{Original lemma}{Original lemmas}
\makeatletter
\newlabel{morphisms-lemma-smooth-etale-over-affine-space}{{054L}{source}{}{sourcepointer-8}{}}
\newlabel{morphisms-lemma-smooth-etale-over-affine-space@cref}{{[originallemma][0][]054L}{[1][0][]source}{}{}{}}
\makeatother
\crefname{originaldefinition}{original definition}{original definitions}
\Crefname{originaldefinition}{Original definition}{Original definitions}
\makeatletter
\newlabel{varieties-definition-relative-frobenius}{{0CC9}{source}{}{sourcepointer-9}{}}
\newlabel{varieties-definition-relative-frobenius@cref}{{[originaldefinition][0][]0CC9}{[1][0][]source}{}{}{}}
\makeatother
\crefname{originallemma}{original lemma}{original lemmas}
\Crefname{originallemma}{Original lemma}{Original lemmas}
\makeatletter
\newlabel{varieties-lemma-relative-frobenius-endomorphism-identity}{{0CCA}{source}{}{sourcepointer-10}{}}
\newlabel{varieties-lemma-relative-frobenius-endomorphism-identity@cref}{{[originallemma][0][]0CCA}{[1][0][]source}{}{}{}}
\makeatother
\crefname{originallemma}{original lemma}{original lemmas}
\Crefname{originallemma}{Original lemma}{Original lemmas}
\makeatletter
\newlabel{varieties-lemma-relative-frobenius-finite}{{0CCD}{source}{}{sourcepointer-11}{}}
\newlabel{varieties-lemma-relative-frobenius-finite@cref}{{[originallemma][0][]0CCD}{[1][0][]source}{}{}{}}
\makeatother
\title{The Garel map of the relative Frobenius of a smooth morphism in characteristic p vanishes except in top degree, where it is surjective}
\author{The Stacks Project Authors}
\date{stacks-project:derham.tex:local-snapshot-20261012; GNU FDL 1.2 or later}
\begin{document}
\maketitle
\noindent\textit{Editorial note.} All delivered mathematical text is the author's own, copied from the pinned Stacks Project chapter file derham.tex; no statement or proof step was rewritten, reordered or removed.\par\medskip
\noindent\textit{Editorial note.} Preamble provenance: the pinned chapter file opens with \textbackslash{}input\{preamble\} and carries no preamble text of its own, so the project's single Stacks-wide preamble staged at sources/stacks-derham/preamble.tex is excerpted verbatim as three spans, lines 8-12 (verbatim package and the project's own reference/slogan/history display environments), lines 14-18 (Xy-pic setup) and lines 154-217 (theorem declarations and mathematical macros). The wrapper class and the amsmath/amssymb/amsthm/graphicx base packages are supplied by the assembler. No macro was invented or redefined.\par\medskip
\noindent\textit{Editorial note.} Same-chapter pointers are located at their \textbackslash{}label lines inside the pinned chapter file; cross-chapter pointers are located in the project's own tag list staged in the snapshot (sources/stacks-chow/tags.txt), at the verbatim lines carrying their canonical tags.\par\medskip
\noindent\textit{Allowed standard prerequisites (editorial).}\begin{itemize}
\item The relative Frobenius and its basic properties. Varieties, definition varieties-definition-relative-frobenius (cited at statement line 4372) and labels varieties-lemma-relative-frobenius-finite (invoked at proof line 4384) and varieties-lemma-relative-frobenius-endomorphism-identity (invoked at proof line 4419). They define \$F\_\{X/S\} : X \textbackslash{}to X\textasciicircum{}\{(p)\}\$, show it is finite, and give the cartesian square expressing the Frobenius of an etale chart over the base.
\item Finiteness, flatness and the local complete intersection property give finite syntomic. More on Morphisms, labels more-morphisms-theorem-criterion-flatness-fibre (invoked at proof line 4388), more-morphisms-lemma-lci-permanence (invoked at proof line 4393) and more-morphisms-lemma-flat-lci (invoked at proof line 4396), and Algebra, label algebra-lemma-CM-over-regular-flat (invoked at proof line 4390). Flatness is checked on fibres through the criterion, the fibre map is flat because the source fibre is Cohen-Macaulay over the regular target, and with the lci property this makes \$F\_\{X/S\}\$ finite syntomic, so that the Garel map \$\textbackslash{}Theta\$ is available.
\item Etale charts of a smooth morphism and the compatibility of \$\textbackslash{}Theta\$ with base change. Morphisms, label morphisms-lemma-smooth-etale-over-affine-space (invoked at proof line 4407), Etale Morphisms, label etale-lemma-relative-frobenius-etale (invoked at proof line 4422) and Same chapter derham.tex, label lemma-etale (invoked at proof line 4425). They provide an etale chart \$U \textbackslash{}to \textbackslash{}mathbf\{A\}\textasciicircum{}d\_S\$ over which the Frobenius square is cartesian with the \$p\$-power map, so that \$\textbackslash{}Theta\$ for \$X\$ is pulled back from \$\textbackslash{}Theta\$ for affine space.
\item The universal and complex computations of the Garel map. Same chapter derham.tex, label example-Garel (invoked at proof line 4439). They allow the formula for \$\textbackslash{}Theta\textasciicircum{}i\$ to be computed in the universal case over \$\textbackslash{}mathbf\{Z\}\$ and, for the needed trace formula, in the complex case after inverting the coordinates.
\end{itemize}\medskip
\ifdefined\hypertarget\hypertarget{sourcepointer-1}{}\fi
\noindent\textit{Source reference.} Original lemma 00R4: sources/stacks-chow/tags.txt, source lines 990--990. Navigation pointer for \textbackslash{}ref\{algebra-lemma-CM-over-regular-flat\} used inside the counted span (line 4390); the argument that consumes it is retained in the span.\par\medskip
\ifdefined\hypertarget\hypertarget{sourcepointer-2}{}\fi
\noindent\textit{Source reference.} Original lemma 0EBS: sources/stacks-chow/tags.txt, source lines 17932--17932. Navigation pointer for \textbackslash{}ref\{etale-lemma-relative-frobenius-etale\} used inside the counted span (line 4422); the argument that consumes it is retained in the span.\par\medskip
\ifdefined\hypertarget\hypertarget{sourcepointer-3}{}\fi
\noindent\textit{Source reference.} Original reference 0FLD: sources/stacks-derham/pinned-original-source.tex, source lines 4308--4308. Navigation pointer for \textbackslash{}ref\{example-Garel\} used inside the counted span (line 4439); the argument that consumes it is retained in the span.\par\medskip
\ifdefined\hypertarget\hypertarget{sourcepointer-4}{}\fi
\noindent\textit{Source reference.} Original lemma 0FLV: sources/stacks-derham/pinned-original-source.tex, source lines 310--310. Navigation pointer for \textbackslash{}ref\{lemma-etale\} used inside the counted span (line 4425); the argument that consumes it is retained in the span.\par\medskip
\ifdefined\hypertarget\hypertarget{sourcepointer-5}{}\fi
\noindent\textit{Source reference.} Original lemma 069K: sources/stacks-chow/tags.txt, source lines 8044--8044. Navigation pointer for \textbackslash{}ref\{more-morphisms-lemma-flat-lci\} used inside the counted span (line 4396); the argument that consumes it is retained in the span.\par\medskip
\ifdefined\hypertarget\hypertarget{sourcepointer-6}{}\fi
\noindent\textit{Source reference.} Original lemma 069M: sources/stacks-chow/tags.txt, source lines 8046--8046. Navigation pointer for \textbackslash{}ref\{more-morphisms-lemma-lci-permanence\} used inside the counted span (line 4393); the argument that consumes it is retained in the span.\par\medskip
\ifdefined\hypertarget\hypertarget{sourcepointer-7}{}\fi
\noindent\textit{Source reference.} Original theorem 039C: sources/stacks-chow/tags.txt, source lines 4267--4267. Navigation pointer for \textbackslash{}ref\{more-morphisms-theorem-criterion-flatness-fibre\} used inside the counted span (line 4388); the argument that consumes it is retained in the span.\par\medskip
\ifdefined\hypertarget\hypertarget{sourcepointer-8}{}\fi
\noindent\textit{Source reference.} Original lemma 054L: sources/stacks-chow/tags.txt, source lines 6644--6644. Navigation pointer for \textbackslash{}ref\{morphisms-lemma-smooth-etale-over-affine-space\} used inside the counted span (line 4407); the argument that consumes it is retained in the span.\par\medskip
\ifdefined\hypertarget\hypertarget{sourcepointer-9}{}\fi
\noindent\textit{Source reference.} Original definition 0CC9: sources/stacks-chow/tags.txt, source lines 15499--15499. Navigation pointer for \textbackslash{}ref\{varieties-definition-relative-frobenius\} used inside the counted span (line 4372); the argument that consumes it is retained in the span.\par\medskip
\ifdefined\hypertarget\hypertarget{sourcepointer-10}{}\fi
\noindent\textit{Source reference.} Original lemma 0CCA: sources/stacks-chow/tags.txt, source lines 15500--15500. Navigation pointer for \textbackslash{}ref\{varieties-lemma-relative-frobenius-endomorphism-identity\} used inside the counted span (line 4419); the argument that consumes it is retained in the span.\par\medskip
\ifdefined\hypertarget\hypertarget{sourcepointer-11}{}\fi
\noindent\textit{Source reference.} Original lemma 0CCD: sources/stacks-chow/tags.txt, source lines 15503--15503. Navigation pointer for \textbackslash{}ref\{varieties-lemma-relative-frobenius-finite\} used inside the counted span (line 4384); the argument that consumes it is retained in the span.\par\medskip
\begin{lemma}
\label{lemma-Garel-map-frobenius-smooth-char-p}
Let $p$ be a prime number. Let $X \to S$ be a smooth morphism
of relative dimension $d$ of schemes in characteristic $p$.
The relative Frobenius $F_{X/S} : X \to X^{(p)}$ of $X/S$
(Varieties, Definition \ref{varieties-definition-relative-frobenius})
is finite syntomic and the corresponding map
$$
\Theta_{X/X^{(p)}} :
F_{X/S, *}\Omega^\bullet_{X/S} \to \Omega^\bullet_{X^{(p)}/S}
$$
is zero in all degrees except in degree $d$ where it defines a
surjection.
\end{lemma}

\begin{proof}
Observe that $F_{X/S}$ is a finite morphism by
Varieties, Lemma \ref{varieties-lemma-relative-frobenius-finite}.
To prove that $F_{X/S}$ is flat, it suffices to show that
the morphism $F_{X/S, s} : X_s \to X^{(p)}_s$ between fibres
is flat for all $s \in S$, see More on Morphisms, Theorem
\ref{more-morphisms-theorem-criterion-flatness-fibre}.
Flatness of $X_s \to X^{(p)}_s$ follows from
Algebra, Lemma \ref{algebra-lemma-CM-over-regular-flat}
(and the finiteness already shown).
By More on Morphisms, Lemma
\ref{more-morphisms-lemma-lci-permanence}
the morphism $F_{X/S}$ is a local complete intersection morphism.
Hence $F_{X/S}$ is finite syntomic (see
More on Morphisms, Lemma \ref{more-morphisms-lemma-flat-lci}).

\medskip\noindent
For every point $x \in X$ we may choose a commutative diagram
$$
\xymatrix{
X \ar[d] & U \ar[l] \ar[d]_\pi \\
S & \mathbf{A}^d_S \ar[l]
}
$$
where $\pi$ is \'etale and $x \in U$ is open in $X$, see
Morphisms, Lemma \ref{morphisms-lemma-smooth-etale-over-affine-space}.
Observe that
$\mathbf{A}^d_S \to \mathbf{A}^d_S$, $(x_1, \ldots, x_d) \mapsto
(x_1^p, \ldots, x_d^p)$ is the relative Frobenius for $\mathcal{A}^d_S$
over $S$. The commutative diagram
$$
\xymatrix{
U \ar[d]_\pi \ar[r]_{F_{X/S}} & U^{(p)} \ar[d]^{\pi^{(p)}} \\
\mathbf{A}^d_S \ar[r]^{x_i \mapsto x_i^p} & \mathbf{A}^d_S
}
$$
of
Varieties, Lemma \ref{varieties-lemma-relative-frobenius-endomorphism-identity}
for $\pi : U \to \mathbf{A}^d_S$ is cartesian by
\'Etale Morphisms, Lemma
\ref{etale-lemma-relative-frobenius-etale}.
Since the construction of $\Theta$ is compatible with base change
and since $\Omega_{U/S} = \pi^*\Omega_{\mathbf{A}^d_S/S}$
(Lemma \ref{lemma-etale})
we conclude that it suffices to show the lemma for
$\mathbf{A}^d_S$.

\medskip\noindent
Let $A$ be a ring of characteristic $p$. Consider the unique $A$-algebra
homomorphism $A[y_1, \ldots, y_d] \to A[x_1, \ldots, x_d]$
sending $y_i$ to $x_i^p$. The arguments above
reduce us to computing the map
$$
\Theta^i : \Omega^i_{A[x_1, \ldots, x_d]/A} \to
\Omega^i_{A[y_1, \ldots, y_d]/A}
$$
We urge the reader to do the computation in this case for themselves.
As in Example \ref{example-Garel} we may reduce this to computing
a formula for $\Theta^i$ in the universal case
$$
\mathbf{Z}[y_1, \ldots, y_d] \to \mathbf{Z}[x_1, \ldots, x_d],\quad
y_i \mapsto x_i^p
$$
In turn, we can find the formula for $\Theta^i$ by computing in the complex
case, i.e., for the $\mathbf{C}$-algebra map
$$
\mathbf{C}[y_1, \ldots, y_d] \to \mathbf{C}[x_1, \ldots, x_d],\quad
y_i \mapsto x_i^p
$$
We may even invert $x_1, \ldots, x_d$ and $y_1, \ldots, y_d$.
In this case, we have $\text{d}x_i = p^{-1} x_i^{- p + 1}\text{d}y_i$.
Hence we see that
\begin{align*}
\Theta^i(
x_1^{e_1} \ldots x_d^{e_d} \text{d}x_1 \wedge \ldots \wedge \text{d}x_i)
& =
p^{-i} \Theta^i(
x_1^{e_1 - p + 1} \ldots x_i^{e_i - p + 1} x_{i + 1}^{e_{i + 1}} \ldots
x_d^{e_d} \text{d}y_1 \wedge \ldots \wedge \text{d}y_i ) \\
& =
p^{-i} \text{Trace}(x_1^{e_1 - p + 1} \ldots x_i^{e_i - p + 1}
x_{i + 1}^{e_{i + 1}} \ldots x_d^{e_d})
\text{d}y_1 \wedge \ldots \wedge \text{d}y_i
\end{align*}
by the properties of $\Theta^i$. An elementary computation shows
that the trace in the expression above is zero unless
$e_1, \ldots, e_i$ are congruent to $-1$ modulo $p$
and $e_{i + 1}, \ldots, e_d$ are divisible by $p$.
Moreover, in this case we obtain
$$
p^{d - i} y_1^{(e_1 - p + 1)/p} \ldots y_i^{(e_i - p + 1)/p}
y_{i + 1}^{e_{i + 1}/p} \ldots y_d^{e_d/p}
\text{d}y_1 \wedge \ldots \wedge \text{d}y_i
$$
We conclude that we get zero in characteristic $p$ unless $d = i$
and in this case we get every possible $d$-form.
\end{proof}

\end{document}
